\begin{theorem}[10 状态同步核的 Mealy 最小性]\label{thm:sync-kernel-mealy-minimality}
沿用 \S\ref{subsec:sync-kernel-counting} 的同步核记号：令 \(K_{\mathrm{sync}}\) 为以 \(Q\) 为状态集、输入字母表为 \(\{0,1,2\}\)、输出字母表为 \(\{0,1\}\) 的确定性 Mealy 转导器（图 \ref{fig:sync_kernel_10_state_graph} 的边标记 \(d/e\) 即给出其迁移与输出）。
则 \(K_{\mathrm{sync}}\) 在可达状态上为 Mealy-最小：其 Nerode 等价类均为单点。
从而任何实现同一顺序函数（同一输入—输出行为）的确定性 Mealy 转导器，其状态数满足 \(\ge 10\)。
\end{theorem}

\begin{proof}
对确定性 Mealy 机的 Nerode 等价关系 \(\sim\)，两个状态 \(q,q'\) 等价当且仅当对所有输入词 \(w\in\{0,1,2\}^\ast\)，从 \(q\) 与 \(q'\) 出发读入 \(w\) 所产生的输出词完全相同。
因此 Mealy 最小性等价于：任意 \(q\neq q'\) 都存在某个区分词 \(w\) 使两者输出不同。
\par
在本核上，可达状态集即 \(Q\)（图 \ref{fig:sync_kernel_10_state_graph}），并且区分词的存在性可被枚举证书化；相应证书的汇总以式 \ref{eq:sync_kernel_10_state_mealy_minimality} 给出（其中 \(|Q/\!\sim|=10\) 表示等价类数目达到状态数本身）。
因此 \(K_{\mathrm{sync}}\) 的每个可达状态均对应互异 Nerode 残余，最小化后状态数仍为 \(10\)，从而结论成立。
\end{proof}

\endinput
